% === §4.1 Objetivo general ===
% Redactado por el Investigador — Fase 3
% Responde directamente a la pregunta de investigación (§2.1)

\noindent\textbf{Desarrollar un ecosistema de agentes autónomos basados en inteligencia artificial que fortalezca el aprendizaje profundo —motivación intrínseca, argumentación técnica y resolución de problemas auténticos— en estudiantes de ingeniería, mediante la integración de un modelo computacional de diagnóstico multidimensional del aprendiz, una arquitectura multiagente pedagógicamente fundada con orquestación basada en grandes modelos de lenguaje y el Model Context Protocol (MCP), y la validación empírica en entornos reales de aula universitaria, alcanzando un nivel de madurez tecnológica TRL~6 y proyectando TRL~7, transferible al entorno educativo y al sector productivo colombiano en un horizonte de 18 meses de ejecución más 2 meses de productos.}

\bigskip

\noindent El objetivo general constituye la respuesta directa a la pregunta de investigación formulada en la sección~\ref{sec:problematica}: articula el \emph{qué} —el ecosistema agéntico como producto tangible de innovación— con el \emph{cómo} —la cadena metodológica que va del diagnóstico del aprendiz a la transferencia tecnológica— y el \emph{para qué} —el fortalecimiento del aprendizaje profundo como impacto verificable en los estudiantes de ingeniería y en el sector productivo colombiano—. La verificación del cumplimiento del objetivo se realizará de forma \textbf{cuantitativa}, a través de tres familias de indicadores: (i) incremento estadísticamente significativo ($p < 0.05$) en las puntuaciones de rúbricas validadas de aprendizaje profundo entre las fases pre y post intervención, con tamaño del efecto reportado ($d$ de Cohen); (ii) métricas de desempeño del ecosistema agéntico en entorno relevante (tasa de finalización de tareas, precisión de la retroalimentación, latencia); y (iii) nivel de madurez tecnológica TRL~6 verificado mediante lista de cotejo estandarizada frente a los criterios de la escala TRL de la NASA/DoD adaptada a \emph{edtech}. Complementariamente, la validación \textbf{cualitativa} se sustentará en el análisis temático de entrevistas semiestructuradas a estudiantes y docentes, la evaluación de la calidad argumentativa en tareas auténticas de ingeniería y el juicio de expertos en educación en ingeniería sobre la pertinencia pedagógica del ecosistema.
